% Generated draft: requires mathematical and source review.
\documentclass[11pt]{article}
\usepackage[margin=25mm]{geometry}
\usepackage{fontspec}
\setmainfont{Latin Modern Roman}
\usepackage{amsmath,amssymb,mathtools,unicode-math}
\setmathfont{Latin Modern Math}
\usepackage{enumitem,needspace,xurl,hyperref}
\setlength{\parindent}{0pt}
\setlength{\parskip}{5pt}
\setlength{\emergencystretch}{3em}
\newcommand{\N}{\mathbb{N}}
\newcommand{\Z}{\mathbb{Z}}
\newcommand{\Q}{\mathbb{Q}}
\newcommand{\R}{\mathbb{R}}
\newcommand{\C}{\mathbb{C}}
\newcommand{\F}{\mathbb{F}}
\newcommand{\A}{\mathbb{A}}
\newcommand{\PP}{\mathbb{P}}
\newcommand{\E}{\mathbb{E}}
\newcommand{\eps}{\varepsilon}
\newcommand{\dd}{\,\mathrm{d}}
\begin{document}
\section*{UnsolvedMath ID: 701682}
\textbf{Source problem code:} catalogue entry
\subsection{Definitions and mathematical statement}
Let
\[
R=\Z[\sqrt{2}]=\{a+b\sqrt{2}:a,b\in\Z\},\qquad
G=\operatorname{SL}_2(R)=\left\{\begin{pmatrix}a&b\\c&d\end{pmatrix}:a,b,c,d\in R,\ ad-bc=1\right\}.
\]
Regard \(G\) as a discrete group. An action of \(G\) on a nonempty set \(X\) is \emph{amenable} if there exists a finitely additive probability measure \(\mu:\mathcal P(X)\to[0,1]\) such that \(\mu(gA)=\mu(A)\) for every \(g\in G\) and every \(A\subseteq X\). The action is \emph{transitive} if, for all \(x,y\in X\), some \(g\in G\) satisfies \(gx=y\), and is \emph{faithful} if the only element fixing every point of \(X\) is the identity.

\textbf{Question.} Does there exist a nonempty set \(X\) with a faithful, transitive, amenable action of \(G\)? Such an \(X\) must be infinite, since \(G\) is infinite and the action is faithful.

A subgroup \(H\leq G\) is \emph{co-amenable} if the left action of \(G\) on \(G/H\) is amenable. The primary question can therefore be expressed as asking whether there exists a co-amenable subgroup \(H\leq G\) with
\[
\operatorname{Core}_G(H)=\bigcap_{g\in G}gHg^{-1}=\{1\}.
\]
The supplied source also states that its question is equivalent to asking whether there exists a co-amenable subgroup \(H\leq G\) with \([G:H]=\infty\). This asserted equivalence is retained as a source qualification: the supplied record does not explain why an infinite-index co-amenable subgroup would yield a faithful transitive amenable action. The infinite-index formulation alone does not explicitly require trivial core.
\subsection{Short English statement}
For \(G=\operatorname{SL}_2(\Z[\sqrt{2}])\), ask whether there is a faithful transitive action on a set that admits a \(G\)-invariant finitely additive probability measure on all its subsets. The source also gives an infinite-index co-amenable subgroup formulation and asserts equivalence, without supplying its justification.
\subsection{Sources}
\begin{itemize}
\item Supplied local UnsolvedMath registry record: permanent ID \(701682\), exact source problem code \texttt{catalogue entry}, title \emph{Kourovka Notebook Problem 20.24}. The registry labels the problem \emph{open}; this is a supplied status, not an independently verified assessment.
\item \emph{The Kourovka Notebook (Unsolved Problems in Group Theory)}, original problem number 20.24, proposed by Y. Glasner and N. Monod in 2022. The supplied background identifies the 20th issue (2022) and the 21st issue edition (2026), arXiv:1401.0300v48. The collection description names E. I. Khukhro and V. D. Mazurov as editors.
\end{itemize}
\end{document}
